\chapter{NỀN TẢNG LÝ THUYẾT VÀ CÔNG NGHỆ SỬ DỤNG}

\section{Nền tảng lý thuyết}

\subsection{Dịch máy thần kinh (Neural Machine Translation)}

Dịch máy thần kinh (NMT) là phương pháp dịch tự động sử dụng mạng nơ-ron học sâu để mô hình hóa
toàn bộ quá trình dịch như một bài toán học máy đầu-đến-cuối (end-to-end). Kiến trúc Transformer
\cite{vaswani2017attention} --- được đề xuất năm 2017 --- đã trở thành nền tảng cho hầu hết các
hệ thống NMT hiện đại, bao gồm Google Neural Machine Translation (GNMT) và các mô hình thương
mại.

Google Cloud Translation~API~v3 sử dụng kiến trúc NMT với các ưu điểm so với phương pháp thống
kê (SMT) truyền thống: độ chính xác cao hơn về ngữ nghĩa, xử lý tốt hơn với câu dài, và hỗ trợ
100+ cặp ngôn ngữ.

\subsection{Glossary trong dịch máy}

Glossary (bảng thuật ngữ) là một danh sách các cặp từ đã biết, trong đó mỗi cặp gồm một từ/cụm
từ nguồn và bản dịch xác định của nó trong ngôn ngữ đích. Khi glossary được tích hợp vào quá
trình dịch, hệ thống dịch sẽ ưu tiên sử dụng bản dịch trong glossary thay vì bản dịch do mô
hình NMT tự sinh ra.

Google Cloud Translation~API~v3 hỗ trợ glossary theo hai loại:
\begin{itemize}
  \item \textit{Unidirectional glossary:} Cặp ngôn ngữ cố định, một chiều.
  \item \textit{Equivalent Term Set glossary:} Hỗ trợ nhiều ngôn ngữ trong một glossary,
    có thể tái sử dụng cho nhiều cặp. Đây là loại được sử dụng trong hệ thống.
\end{itemize}

Glossary được lưu dưới dạng tệp CSV trên Google Cloud Storage và đăng ký với Google Translation
API. Mỗi khi thuật ngữ trong thư viện chung thay đổi, hệ thống tự động tạo lại tệp CSV và cập
nhật glossary tương ứng qua background thread.

\subsection{Nhận dạng ký tự quang học (OCR)}

OCR (Optical Character Recognition) là công nghệ nhận diện văn bản từ hình ảnh hoặc tài liệu
scan. Trong hệ thống, Azure Document Intelligence~--- dịch vụ OCR của Microsoft Azure~--- được
sử dụng để trích xuất nội dung từ các tệp PDF dạng scan trước khi dịch.

Azure Document Intelligence cung cấp khả năng nhận diện bố cục trang (trang, bảng, hình ảnh),
hỗ trợ đa ngôn ngữ Á-Âu và độ chính xác cao hơn so với Tesseract OCR truyền thống với tài
liệu kỹ thuật có bảng biểu và ký tự đặc biệt.

\section{Công nghệ sử dụng}

\subsection{Frontend}

\begin{table}[H]
\centering
\caption{Công nghệ Frontend}
\label{tab:tech_fe}
\renewcommand{\arraystretch}{1.3}
\begin{tabularx}{\textwidth}{|l|l|X|}
\hline
\textbf{Công nghệ} & \textbf{Phiên bản} & \textbf{Mục đích} \\
\hline
React              & 18.3              & Thư viện UI, xây dựng SPA theo mô hình component \\
\hline
Vite               & 6.x               & Build tool và dev server; Hot Module Replacement nhanh \\
\hline
React Router DOM   & 7.x               & Định tuyến phía client \\
\hline
Tailwind CSS       & 4.x               & Utility-first CSS framework \\
\hline
Ant Design (antd)  & 5.x               & UI component library: table, modal, form, notification \\
\hline
Axios              & 1.6+              & HTTP client gọi REST~API \\
\hline
react-toastify     & --                & Hiển thị thông báo toast \\
\hline
MSAL Browser       & --                & Microsoft Authentication Library cho Azure AD SSO \\
\hline
jwt-decode         & --                & Giải mã JWT token phía client \\
\hline
@bytescale/upload-widget & --         & Widget tải tệp lên S3 \\
\hline
recharts           & --                & Biểu đồ thống kê (keyword analytics) \\
\hline
\end{tabularx}
\end{table}

React được chọn vì kiến trúc component-based phù hợp với giao diện có nhiều widget tương tác (upload, language selector, keyword table). Vite được chọn thay Webpack vì thời gian build và hot-reload nhanh hơn đáng kể trong quá trình phát triển.

\subsection{Backend}

\begin{table}[H]
\centering
\caption{Công nghệ Backend}
\label{tab:tech_be}
\renewcommand{\arraystretch}{1.3}
\begin{tabularx}{\textwidth}{|l|l|X|}
\hline
\textbf{Công nghệ}           & \textbf{Phiên bản} & \textbf{Mục đích} \\
\hline
Python                       & 3.9+              & Ngôn ngữ lập trình chính \\
\hline
Django                       & 6.0               & Web framework, ORM, migration \\
\hline
Django REST Framework (DRF)  & 3.16.1            & Xây dựng REST~API \\
\hline
djangorestframework-simplejwt & 5.5.1            & JWT authentication (dev) \\
\hline
psycopg2-binary              & --                & Kết nối PostgreSQL \\
\hline
python-dotenv                & --                & Quản lý biến môi trường \\
\hline
\end{tabularx}
\end{table}

Django được chọn so với FastAPI và Flask vì: tích hợp ORM và migration mạnh mẽ, DRF cung cấp
serializer và permission class hoàn chỉnh, hệ thống admin tích hợp sẵn.

\subsection{Dịch thuật và xử lý tài liệu}

\begin{table}[H]
\centering
\caption{Thư viện xử lý tài liệu và dịch thuật}
\label{tab:tech_translate}
\renewcommand{\arraystretch}{1.3}
\begin{tabularx}{\textwidth}{|l|X|}
\hline
\textbf{Thư viện / Dịch vụ}   & \textbf{Mục đích} \\
\hline
google-cloud-translate v3     & API dịch thuật chính, hỗ trợ custom glossary \\
\hline
google-cloud-storage          & Upload/download CSV glossary lên GCS \\
\hline
azure-ai-documentintelligence & OCR tài liệu PDF scan \\
\hline
convertapi                    & Chuyển đổi PDF→DOCX và XLS/CSV→XLSX \\
\hline
python-docx                   & Đọc/ghi tệp DOCX; trích xuất text từ paragraphs và tables \\
\hline
lxml                          & Parse và ghi XML trong DOCX (zip-based); merge run có cùng format \\
\hline
openpyxl                      & Đọc/ghi tệp XLSX; dịch từng cell \\
\hline
python-pptx                   & Đọc/ghi tệp PPTX; dịch text trong shapes và slides \\
\hline
pypdf / PyMuPDF               & Trích xuất text và metadata từ PDF để phân loại \\
\hline
Pillow / opencv-python        & Xử lý ảnh trong pipeline OCR \\
\hline
boto3                         & Tương tác với AWS~S3 \\
\hline
\end{tabularx}
\end{table}

\textbf{Lý do chọn Google Cloud Translation~API~v3 thay vì v2:}
API~v3 hỗ trợ custom glossary với Equivalent Term Set (nhiều ngôn ngữ một glossary), batch
translation (dịch nhiều chuỗi trong một lần gọi API), và model selection. Tính năng glossary là
yêu cầu bắt buộc của hệ thống nên v3 là lựa chọn duy nhất phù hợp.

\subsection{Hạ tầng và triển khai}

\begin{table}[H]
\centering
\caption{Công nghệ hạ tầng và triển khai}
\label{tab:tech_infra}
\renewcommand{\arraystretch}{1.3}
\begin{tabularx}{\textwidth}{|l|X|}
\hline
\textbf{Công nghệ}         & \textbf{Mục đích} \\
\hline
AWS S3                     & Lưu trữ tệp gốc và tệp dịch với UUID-based key \\
\hline
Amazon Cognito             & SSO doanh nghiệp qua OIDC/OAuth2; ALB tích hợp sẵn \\
\hline
AWS ALB                    & Cân bằng tải; inject header \texttt{x-amzn-oidc-data} cho xác thực \\
\hline
PostgreSQL 12+             & Cơ sở dữ liệu quan hệ (production) \\
\hline
Docker                     & Container hóa backend và frontend; đảm bảo tái tạo môi trường \\
\hline
Nginx                      & Web server phục vụ frontend SPA; reverse proxy trong container \\
\hline
Google Cloud Storage       & Lưu trữ tệp CSV glossary cho Google Translation API \\
\hline
\end{tabularx}
\end{table}

\subsection{Cơ sở dữ liệu}

PostgreSQL được chọn là cơ sở dữ liệu quan hệ chính vì: hỗ trợ \texttt{JSONField} (dùng cho
\texttt{keyword\_details} trong Notification), indexing mạnh mẽ, tương thích tốt với Django ORM.
SQLite được dùng như fallback trong môi trường phát triển để đơn giản hóa setup.
